\chapter{Основен резултат}

Нека $F:\R^n\times [0,T]\rightrightarrows \R^n$ има непразни затворени образи и $F(x,t)$ е измеримо по $t$ и за всяко $t\in [0,T]$ $F(\cdot, t)$ е $l(t)$-липшицова за някое $l\in L^1([0,T])$.
Нека $\Point\in\Univ$.

Първо да докажем някои по-технически свойства, свързани с измеримост:
\begin{lemma}
    Нека $(v,w)\in X\times Y$. Тогава са измерими следните:
    \begin{enumerate}
        \item изображението $t \mapsto \Gr$
        \item за всяко $\delta > 0$ изображението $G_{\delta}(t) := \bar{\Ball}_{\delta}(\Pointt)\cap\Gr$
        \item за всеки $r,s > 0$ множеството
            \begin{align*}
                D_{r,s} = \{t\in [0,T]\mid\ \forall z\in\bar{\Ball}_{r}(\Pointt)&\cap\Gr\ \forall \tau\in(0,s)\\
                  &[z + \tau \bar{\Ball}_{\varepsilon}(\Vect)]\cap\Gr\neq\emptyset\}.
            \end{align*}
    \end{enumerate}
\end{lemma}
\begin{proof}
    1) За всяко $x\in\R^n$ знаем, че $F(x,\cdot)$ е измеримо. Нека вземем функциите $f_n^x(t)$, т.ч. $F(x,t) = \overline{\{f^x_n(t) \mid n\in\N\}}$. Тогава
    \[
        \Gr = \overline{\{(q,f_n^q(t)) \mid n\in\N\ \&\ q\in\Q^n\}}
    \]
    Наистина ако $(x,y)\in\Gr$, поради непрекъснатостта на $F(\cdot, t)$, за всяко $n\geq 1$ съществува $\delta_n \in (0, \frac{1}{n})$, т.ч. $F(x,t) \subseteq F(x', t) + \frac{1}{n}\bar{\Ball}$ за всяко $x'\in \Ball_{\delta_n}(x)$.
    Да изберем $q_n\in\Q^n$ т.ч. $\lVert q_n - x\rVert \leq \delta_n$ и $y_n\in F(q_n, t)$, т.ч. $\lVert y - y_n\rVert \leq \frac{1}{n}$.
    Поради гъстотата на $f_i^{q_n}(t)$ в образа, можем да намерим $i_n$, т.ч. $\lVert f_{i_n}^{q_n}(t) - y_n \rVert \leq \frac{1}{n}$.
    Тогава $(q_n, f_{i_n}^{q_n}(t)) \xrightarrow[n\to\infty]{} (x,y)$.
    % \[
    %     \lVert (q_n, f_{i_n}^{q_n}(t)) - (x,y)\rVert =
    %     % \lVert q_n - x\rVert + \lVert f_{i_n}^{q_n} - y \rVert \leq
    %     \lVert q_n - x\rVert + \lVert f_{i_n}^{q_n}(t) - y_n \rVert + \lVert y_n - y \rVert \leq
    %     \frac{3}{n}
    %     \xrightarrow[n\to\infty]{} 0
    % \]
    Обратното включване следва директно от Теорема \ref{ClosedGraph}.\\
    2) Теорема $\ref{MeasurableBall}$ ни дава измеримостта на $\bar{\Ball}_{\delta}(\Pointt)$, а Теорема $\ref{MeasurableCap}$ - измеримостта на сечението.\\
    3) Да забележим, че условието
    \[
        [z + \tau \bar{\Ball}_{\varepsilon}(\Vect)]\cap\Gr\neq\emptyset
    \]
    е еквивалентно на
    \[
        f(z, \tau, t) := \frac{d(z+\tau\Vect, \Gr)}{\tau} \leq \varepsilon
    \]
    където $f(z,\tau,t)$ е непрекъсната по $z$ и $\tau$ и измерима по $t$.
    Тогава
    \[
        D_{r,s} = \{t\in [0,T] \mid g(t) := \sup_{\substack{z\in G_r(t)\\\tau\in(0,s)}} f(z,\tau,t)\leq\varepsilon\}.
    \]
    Трябва само да се уверим, че $g$ е измерима.

    Нека $h_n(t)$ са измерими функции т.ч. $G_r(t) = \overline{\{h_n(t) \mid n\in\N\}}$. Тогава
    \[
        \{ g(t) > \alpha \} = \bigcup_{\substack{n\in\N\\\tau\in(0,s)}} \{ f(h_n(t),\tau,t) > \alpha \}.
    \]
    Наистина, ако $g(t) > \alpha$, тогава съществуват $z_0\in G_r(t)$ и $\tau_0\in (0,s)$, т.ч. $f(z_0, \tau_0, t) > \alpha$.
    Понеже $f$ е непрекъсната по $(z,\tau)$, то съществува $U$ - околност на $(z_0, \tau_0)$ в $G_r(t)\times (0,s)$, т.ч. $f(z,\tau,t) > \alpha$ за всяко $(z,\tau)\in U$.
    Поради гъстотата можем да намерим $n\in\N$ и $\tau\in\Q$, т.ч. $(h_n(t), \tau)\in U$ и тогава $f(h_n(t), \tau, t) > \alpha$.
    Другото включване е ясно.
\end{proof}

\begin{lemma} \label{L2}
    Нека $(v,w)\in\Univ$, $E\subseteq [0,T]$ е измеримо и $\nu > 0$. Тогава съществува компакт $E'\subseteq E$ с $\mu(E\setminus E') < \nu$, т.ч. са измерими:
    \begin{enumerate}
        \item за всеки $\delta,\lambda,\varepsilon > 0$ изображенията $H_{\delta,\lambda,\varepsilon}:E'\rightrightarrows \R^{2n}$ зададени с
            \[
                H_{\delta,\lambda,\varepsilon}(t) := \{z\in G_{\delta}(t) \mid [z + (0,\lambda)\Ball_{\varepsilon}(\Vect)]\cap\Gr = \emptyset\},
            \]
        \item за всяко $\delta,\lambda,\varepsilon>0$ множеството
            \begin{align*}
                L_{\delta,\lambda,\varepsilon}=\{t\in E' \mid\ &\exists z\in \clBall_{\delta}(\Pointt)\cap\Gr\\
                   &[z + (0, \lambda)\Ball_{\varepsilon}(\Vect)\cap\Gr] = \emptyset\},
            \end{align*}
        \item за всяко $s>0$ множеството
            \begin{align*}
                E_s = \{t\in E' \mid \forall \delta > 0&\ \exists z \in \clBall_{\delta}(\Pointt)\cap\Gr\ \exists \lambda > 0\\
                 &[z + (0, \lambda)\Ball_s(\Vect)]\cap\Gr = \emptyset\}.
                \end{align*}
    \end{enumerate}
    Също така изображенията $H_{\delta,\lambda,\varepsilon}$ имат затворени образи върху $E'$.
\end{lemma}
\begin{proof}
    Да забележим първо, че условието
    \[
        [z+(0, \lambda)\Ball_{\varepsilon}(\Vect)]\cap\Gr=\emptyset
    \]
    е еквивалентно на
    \[
        g(t, z):=\inf_{\tau\in(0,\lambda)}\frac{d(z+\tau\Vect, \Gr)}{\tau} \geq \varepsilon.
    \]
    $g(t,z)$ е измерима по $t$ и полунепрекъсната отгоре по $z$.
    Измеримостта се показва аналогично на горната теорема.
    За полунепрекъснатостта е достатъчно да забележим, че ако $f(x,y)$ е непрекъсната по $x$,\\ $g(x) = \inf_{y\in A} f(x,y)$ и $\varepsilon > 0$, имаме
    \[
        g(x) + \varepsilon > f(x,y_0) = \limsup_{x'\to x} f(x',y_0) \geq \limsup_{x'\to x} g(x')
    \]
    Оставяйки $\varepsilon\to 0$, получаваме искания резултат.
    Тогава според теорема \ref{Scorza} съществува компакт $E'\subseteq E$ с $\mu(E\setminus E')<\nu$, т.ч. функцията $g$ е измерима върху $E'\times\R^{2n}$.
    \begin{enumerate}
        \item
            Графиката на $H_{\delta,\lambda,\varepsilon}$ можем да представим като
            \[
                Gr H_{\delta,\lambda,\varepsilon} = \{(t,z) \mid (t,z)\in Gr G_{\delta}\ \&\ g(t,z)\geq\varepsilon\} =
                Gr G_{\delta} \cap g\vert_{E'\times\R^{2n}}^{-1}([\varepsilon, \infty)),
            \]
            което е $\mathcal{L}\otimes\mathcal{B}$-измеримо.
            Затвореността е ясна от представянето $H_{\delta,\lambda,\varepsilon}(t) = G_\delta(t) \cap g(t,\cdot)^{-1}([\varepsilon,\infty))$ на $H_{\delta,\lambda,\varepsilon}(t)$ и фактът, че $g$ е полунепрекъсната отгоре по $z$.
            % Сега да видим затвореността.\\
            % Нека $z_n\in H_{\delta,\lambda,\varepsilon}(t) = G_\delta(t) \cap \{z \mid g(t,z) \geq\varepsilon\}$ и $z_n\to z_0$.
            % Понеже $G_\delta(t)$ е затворено, $z_0\in G_\delta(t)$. Също така за всяко $n\in\N$, имаме $g(t,z_n)\geq\varepsilon$. От полунепрекъснатостта отгоре получаваме
            % \[
            %     g(t,z_0) \geq \limsup_{n\to\infty} g(t,z_n) \geq\varepsilon.
            % \]
            % Така $z_0\in H_{\delta,\lambda,\varepsilon}(t)$.
        \item
            Имаме следното представяне
            \[
                L_{\delta,\lambda,\varepsilon} =
                \{t\in E' \mid \exists z\in G_{\delta}(t)\ g(t,z)\geq\varepsilon\}=
                dom(H_{\delta,\lambda,\varepsilon}),
            \]
            където $dom(H_{\delta,\lambda,\varepsilon}) = H_{\delta,\lambda,\varepsilon}^{-1}(\R^{2n})$ са измерими.
        \item
            Имаме следното представяне
            \[
                E_s =
                \bigcap_{\delta\in\Q^+} \bigcup_{\lambda\in\Q^+} \{t\in E' \mid \exists z\in G_{\delta}(t)\ g(t,z)\geq s\} =
                \bigcap_{\delta\in\Q^+} \bigcup_{\lambda\in\Q^+} L_{\delta,\lambda,s}
            \]
    \end{enumerate}
\end{proof}

\begin{proposition}
    За всяко $(\bar{x},\bar{y})\in B$
    \small
    \begin{displaymath}
        \Tangent_B(\Point) \subseteq \{ \Vec\in\Univ \mid \Vect\in T_{\Gr}(\Pointt) \text{ п.н. в } [0,T]\}.
    \end{displaymath}
\end{proposition}

\begin{proof}
    % фаза 1: равномерен епсилон
    Да допуснем, че $(v,w)\in\Univ$ не е в дясното множество.
    Тогава според теорема \ref{Treiman} съществува $E\subseteq [0,T]$ с положителна мярка, т.ч. за всяко $t\in E$ съществува $\varepsilon_t > 0$, т.ч. за всяко $\delta > 0$ съществуват $(x,y)\in \clBall_\delta(\Pointt)\cap\Gr$ и $\lambda > 0$, т.ч.
    \[
        [(x,y) + (0, \lambda)\Ball_{\varepsilon_t}(\Vect)]\cap\Gr = \emptyset
    \]
    След евентуално свиване на $E$ (с $\nu = \frac{\mu(E)}{2}$) можем да предположим, че в $E$ измеримостите от теорема \ref{L2} са в сила.

    Да разгледаме множествата
    \begin{align*}
        E_s = \{t\in E \mid \forall \delta > 0&\ \exists (x,y) \in \bar{\Ball}_\delta(\Pointt)\cap\Gr\ \exists \lambda > 0\\
         &[(x,y) + (0, \lambda)\Ball_s(\Vect)]\cap\Gr = \emptyset\}.
    \end{align*}
    Забелязваме, че
    \[
        E_s \subseteq E_t,\ \text{когато}\ t\leq s,\ \text{както и}\ 
        E = \bigcup_{n = 1}^{\infty} E_{\frac{1}{n}}
    \]
    и тогава
    \[
        0 < \mu(E) = \mu\left(\bigcup_{n=1}^{\infty} E_{\frac{1}{n}}\right) =
        \lim_{n\to\infty} \mu(E_{\frac{1}{n}})
    \]
    и значи съществува някое $\varepsilon_0 := \frac{1}{n_0} > 0$, т.ч. $\mu(\bar{E}:=E_{\varepsilon_0}) > 0$.\\
    Така за всяко $t\in \bar{E}$ и $\delta > 0$ имаме:
    \begin{align*}
        \exists z \in \clBall_{\delta}(\Pointt)\cap&\Gr\ \exists \lambda > 0\\
        &[z + (0, \lambda)\Ball_{\varepsilon_0}(\Vect)]\cap\Gr = \emptyset.
    \end{align*}

    % даваме епсилон и взимаме делта
    Нека $\varepsilon'=\frac{\varepsilon_0}{2}\min\{L,\frac{\mu(\bar{E})}{4}\}$ и $\delta' > 0$ е произволно. Полагаме $\delta := \frac{\delta'}{1+T}$.

    % фаза 2: равномерни ламбда и z
    Дефинираме
    \begin{align*}
        \Lambda_s=\{&t\in\bar{E} \mid \exists z\in \clBall_{\delta}(\Pointt)\cap\Gr\\
        &[z + (0, s)\Ball_{\varepsilon_0}(\Vect)]\cap\Gr = \emptyset\} =
        \bar{E} \cap L_{\delta,s,\varepsilon_0}.
    \end{align*}
    Отново
    \[
        \Lambda_s \subseteq \Lambda_t,\ \text{когато}\ t\leq s,\ \text{както и}\ 
        \bar{E} = \bigcup_{n = 1}^{\infty} \Lambda_{\frac{1}{n}}
    \]
    и значи
    \[
        \mu(\bar{E}) =
        \mu\left(\bigcup_{n=1}^{\infty} \Lambda_{\frac{1}{n}}\right) =
        \lim_{n\to\infty} \mu(\Lambda_{\frac{1}{n}})
    \]
    и значи има някое $\lambda' := \frac{1}{n_0}$, т.ч.
    \[
        \mu(\Set := \Lambda_\frac{1}{n_0}) > \frac{\mu(\bar{E})}{2}.
    \]
    За така построените $\Set$ и $\lambda'$ е вярно, че за всяко $t\in\Set$ съществува $z\in \clBall_{\delta}(\Pointt)\cap\Gr$, т.ч.
    \[
        [z + (0, \lambda')\Ball_{\varepsilon_0}(\Vect)]\cap\Gr = \emptyset
    \]
    Така изображението $H_{\delta,\lambda',\varepsilon_0}:\Set\to\R^{2n}$ има непразни затворени образи. Взимаме измерима селекция $z(t) = (x(t),y(t)) \in H_{\delta,\lambda',\varepsilon_0}(t)$.\\
    Ще използваме теорема \ref{Treiman} за да покажем, че $\Vec\notin \Tangent_B(\Point)$.

    % фаза 3: довършваме
    Дефинираме $(\tilde{x}, \tilde{y}) \in\Univ$ като
    \[
        (\tilde{x}(t), \tilde{y}(t))=
        \begin{cases}
            (x(t), y(t)), &\text{ ако } t\in \Set \\
            \Pointt, &\text{ ако } t\notin \Set \\
        \end{cases}
    \]
    Тогава
    \[
        (\tilde{x}(t), \tilde{y}(t))\in\Gr,
    \]
    \[
        \lVert (\tilde{x}, \tilde{y}) - \Point \rVert =
        L \lVert \tilde{x} - \bar{x} \rVert_\infty + \lVert \tilde{y} - \bar{y} \rVert_1 =
    \]
    \[
        \sup_{t \in\Set}L\lVert \tilde{x}(t) - \bar{x}(t) \rVert + \int_{\Set}\lVert \tilde{y}(t) - \bar{y}(t) \rVert dt\leq
    \]
    \[
        \sup_{t \in\Set}\lVert (\tilde{x}(t), \tilde{y}(t)) - \Pointt \rVert + \int_{\Set}\lVert (\tilde{x}(t), \tilde{y}(t)) - \Pointt \rVert dt\leq
    \]
    \[
        \delta + \int_{\Set}\delta dt \leq (1+T)\delta = \delta'
    \]
    Ще докажем, че
    \[
        (\tilde{x}, \tilde{y}) + (0,\lambda')(\Ball_{\varepsilon'}(\Vec)) \cap B = \emptyset
    \]
    Нека $(p,q)\in \Ball_{\varepsilon'}(\Vec)$ е произволно. Това значи, че
    \[
        \lVert p - v\rVert_\infty \leq \frac{1}{L} \lVert (p,q) - \Vec\rVert <
        \frac{\varepsilon'}{L} \leq \frac{\varepsilon_0}{2}
    \]
    както и, че ако разгледаме $C=\{t\in\Set \mid \lVert q(t) - w(t)\rVert \geq \frac{\varepsilon_0}{2}\}$, имаме
    \[
        \frac{\varepsilon_0}{2} \mu(C) = \int_C \frac{\varepsilon_0}{2}dt \leq
        \int_C \lVert q(t) - w(t) \rVert dt \leq \lVert q - w \rVert_1 \leq \lVert (p,q)-\Vec \rVert <
        \varepsilon'
    \]
    значи $\mu(C) < \frac{2 \varepsilon'}{\varepsilon_0}\leq \frac{\mu(\bar{E})}{4} < \frac{\mu(\Set)}{2}$.
    откъдето $\mu(\Set\setminus C)>0$.\\
    За всяко $t\in\Set\setminus C$ имаме
    \[
        \lVert p(t) - v(t) \rVert \leq \frac{\varepsilon_0}{2}
        \quad \text{и} \quad
        \lVert q(t) - w(t) \rVert < \frac{\varepsilon_0}{2}
    \]
    и значи $\lVert (p(t),q(t)) - \Vect\rVert < \varepsilon_0$.
    Но за $t\in\Set$ ние знаем, че
    \[
        [z(t) + (0, \lambda')\Ball_{\varepsilon_0}(\Vect)]\cap\Gr = \emptyset
    \]
    откъдето $(\tilde{x}, \tilde{y}) + (0,\lambda')(p,q)\cap B = \emptyset$. С това доказахме теоремата.
\end{proof}

Да забележим, че в доказателството на тази теорема ползвахме само непрекъснатостта по $x$. Липшицовостта на изображението ни трябва само за другата посока:

\begin{proposition}
    За всяко $(\bar{x},\bar{y})\in B$
    \small
    \begin{displaymath}
        \Tangent_B(\Point) \supseteq \{ \Vec\in\Univ \mid \Vect\in T_{\Gr}(\Pointt) \text{ п.н. в } [0,T]\}.
    \end{displaymath}
\end{proposition}
\begin{proof}
    Да вземем произволно $\Vec$ от дясната страна.
    Ще ползваме теорема \ref{Geometric} за да покажем, че $\Vec\in \Tangent_B(\Point)$. Нека $\varepsilon > 0$.

    Понеже $l(t)$ и $\lVert w(t)\rVert$ са сумируеми, това означава, че съществува $\nu > 0$, т.ч.
    \[
        \int_S l(t) dt < \frac{\varepsilon}{3(\lVert v \rVert_{\infty} + 1)}
        \text{ и }
        \int_S \lVert w(t)\rVert dt < \frac{\varepsilon}{3}
        \text{ когато }
        \mu(S)\leq\nu
    \]

    Имаме множество $E$ с $\mu(E)=T$, т.ч. за всяко $t\in E$, $\Vect\in \Tangent_{\Gr}(\Pointt)$.
    Нека $\varepsilon' = \frac{\varepsilon}{3(L+T)}$, тогава от теорема \ref{Geometric} за всяко $t\in E$ съществуват $\delta > 0$ и $\lambda > 0$, т.ч. за всеки $(x,y)\in \bar{\Ball}_{\delta}(\Pointt)$ и $\tau\in(0,\lambda)$ имаме
    \[
        [(x,y)+ \tau \clBall_{\varepsilon'}(\Vect)]\cap\Gr\neq\emptyset
    \]
    Да разгледаме множествата
    \begin{align*}
        \Delta_{r,s} = \{t\in E\mid\ &\forall (x,y)\in\bar{\Ball}_{r}(\Pointt)\cap\Gr\ \forall \tau\in(0,s)\\
        &[(x,y) + \tau \bar{\Ball}_{\varepsilon'}(\Vect)]\cap\Gr\neq\emptyset\}=
        D_{r,s}\cap E.
    \end{align*}
    Тогава
    \[
        \Delta_{r,s} \subseteq \Delta_{r,t} \text{ за всяко } r \text{ и всяко } s \leq t
    \]
    \[
        \Delta_{r,s} \subseteq \Delta_{t,s} \text{ за всяко } s \text{ и всяко } r \leq t
    \]
    \[
        E = \bigcup_{n = 1}^{\infty} \bigcup_{m = 1}^{\infty} \Delta_{\frac{1}{n},\frac{1}{m}}
    \]
    откъдето
    \[
        T = \mu(E) =
        \mu\left(\bigcup_{n = 1}^{\infty} \bigcup_{m = 1}^{\infty} \Delta_{\frac{1}{n},\frac{1}{m}}\right) =
        \lim_{n\to\infty}\lim_{m\to\infty} \mu(\Delta_{\frac{1}{n},\frac{1}{m}})
    \]
    и значи съществуват $\delta_0 := \frac{1}{n_0}$ и $\lambda_0 := \frac{1}{m_0}$, т.ч.
    \[
        \mu(\bar{E} := \Delta_{\delta_0,\lambda_0}) \geq T - \frac{\nu}{2}.
    \]
    Взимаме $\delta = \min\{L, \frac{\nu}{2}\}\frac{\delta_0}{2}$ и $\lambda = \lambda_0$.
    Сега нека $(x,y)\in \bar{\Ball}_{\delta}(\Point)$ и $\tau\in(0,\lambda)$ са произволни.

    Да разгледаме множеството $E_0 = \{t\in E \mid \lVert y(t) - \bar{y}(t)\rVert > \frac{\delta_0}{2}\}$. За него имаме
    \[
        \delta \geq
        \lVert y - \bar{y}\rVert_1 =
        % \int_0^T \lVert y(t) - \bar{y}(t) \rVert dt \geq
        \int_{E_0} \lVert y(t) - \bar{y}(t) \rVert dt \geq
        \int_{E_0} \frac{\delta_0}{2} dt =
        \frac{\delta_0}{2} \mu(E_0)
    \]
    Откъдето $\mu(E_0) \leq \frac{2\delta}{\delta_0} \leq \frac{\nu}{2}$.

    Разбиваме $[0,T]$ на $\Set:=\bar{E}\setminus E_0$ и допълнението му - $\Set^C=\bar{E}^C \cup E_0$.\\
    За всяко $t\in\Set$ имаме $\lVert y(t) - \bar{y}(t)\rVert \leq \frac{\delta_0}{2}$ както и
    \[
        L\frac{\delta_0}{2} \geq \delta \geq L\lVert x - \bar{x}\rVert_{\infty} \geq L\lVert x(t) - \bar{x}(t)\rVert
    \]
    и значи $(x(t), y(t))\in \bar{\Ball}_{\delta_0}(\Pointt)$. Но това значи, че
    \[
        [(x(t),y(t))+ \tau \clBall_{\varepsilon'}(\Vect)]\cap\Gr\neq\emptyset
    \]
    За допълнението на $\Set$ имаме
    \[
        \mu(\Set^C) \leq \mu(\bar{E}^C) + \mu(E_0) \leq
        \frac{\nu}{2} + \frac{\nu}{2} = \nu.
    \]
    Там ще ползваме малко по-грубата оценка:
    \[
        d(y(t) + \tau w(t), F(x(t) + \tau v(t), t)) \leq
    \]
    \[
        d(y(t) + \tau w(t), y(t)) + d(y(t), F(x(t), t)) + d_H(F(x(t), t), F(x(t) + \tau v(t), t)) \leq
    \]
    \[
        \tau \lVert w(t)\rVert + 0 + l(t)\tau \lVert v(t)\rVert <
        \tau (\lVert w(t)\rVert + l(t)(\lVert v(t)\rVert + 1)).
    \]
    Да положим $\mu(t) := \lVert w(t)\rVert + l(t)(\lVert v(t)\rVert + 1)$. Горната оценка ни дава
    \[
        [x(t) + \tau v(t), y(t) + \tau \mu(t)\clBall(w(t))] \cap\Gr\neq\emptyset
    \]
    Строим изображението
    \[
        H(t) =
        \begin{cases}
            [(x(t),y(t))+ \tau \bar{\Ball}_{\varepsilon'}(\Vect)]\cap\Gr, &\text{ ако } t\in\Set\\
            [x(t) + \tau v(t), y(t) + \tau \mu(t) \bar{\Ball}(w(t))] \cap\Gr, &\text{ ако } t\in\Set^C
        \end{cases}
        .
    \]
    От горните съображения $H$ е измеримо с непразни затворени образи. Тогава можем да изберем измерима селекция $z(t)\in H(t)$. Избираме изображенията $p$ и $q$, така че да изпълняват $z(t) = (x(t) + \tau p(t), y(t) + \tau q(t))$.
    За така избраните изображения е ясно, че $(x,y) + \tau (p,q)\in B$. Остава да показжем само, че $(p,q)\in \bar{\Ball}_{\varepsilon}(\Vec)$. Пресмятаме:
    \[
        \lVert (p, q) - (v, w)\rVert =
        L \lVert p - v \rVert_{\infty} + \lVert q - w\rVert_1 =
    \]
    \[
        L \sup_{x\in\Set} \lVert p(t) - v(t)\rVert + \int_{\Set} \lVert q(t) - w(t)\rVert dt + \int_{\Set^C} \lVert q(t) - w(t)\rVert dt \leq
    \]
    \[
        L \varepsilon' + \int_{\Set} \varepsilon' dt + \int_{\Set^C} \mu(t) dt \leq
        (L + T) \varepsilon' + \int_{\Set^C} \lVert w(t)\rVert + l(t) (\lVert v(t)\rVert + 1) dt \leq
    \]
    \[
        \frac{\varepsilon}{3} + \int_{\Set^C} \lVert w(t)\rVert dt + (\lVert v(t)\rVert_{\infty} + 1) \int_{\Set^C} l(t) dt \leq
        \frac{\varepsilon}{3} + \frac{\varepsilon}{3} + \frac{\varepsilon}{3} = \varepsilon
    \]
\end{proof}

\noindent И така доказахме

\begin{theorem} \label{BigCone}
    \small
    \begin{displaymath}
        \Tangent_B(\Point) = \{ \Vec\in\Univ \mid \Vect\in T_{\Gr}(\Pointt) \text{ п.н. в } [0,T]\}.
    \end{displaymath}
\end{theorem}

Остана само да докажем

\begin{theorem}
    Нека $\bar{x}$ е решение на (1). Тогава $A$ и $B$ са субтрансверзални в точката $(\bar{x},\dot{\bar{x}})\in A\cap B$.
\end{theorem}
\begin{proof}
    Ще използваме теорема \ref{Subsiff} с $\delta = 1$.\\
    Нека $(x,\dot{x})\in A\cap B((\bar{x},\dot{\bar{x}}))$.
    Означаме $p(t) := d(\dot{x}(t), F(x(t), t))$.
    Да забележим, че за всеки $(x_1,y_1)\in B$ имаме:
    \[
        p(t) \leq
        d(\dot{x}(t), y_1(t)) + d(y_1(t), F(x_1(t),t)) + d_H(F(x_1(t), t), F(x(t), t)) \leq
    \]
    \[
        \lVert \dot{x}(t) - y_1(t)\rVert + 0 + l(t) \lVert x_1(t) - x(t)\rVert
    \]
    откъдето
    \[
        \int_0^T p(t) dt \leq
        \int_0^T \lVert \dot{x}(t) - y_1(t)\rVert + l(t) \lVert x_1(t) - x(t)\rVert dt \leq
    \]
    \[
        \lVert \dot{x} - y_1 \rVert_1 + \lVert x_1 - x\rVert_\infty \int_0^T l(t)dt =
        \lVert \dot{x} - y_1 \rVert_1 + L\lVert x_1 - x\rVert_\infty =
        \lVert (x_1, y_1) - (x,\dot{x})\rVert
    \]
    понеже $(x_1,y_1)\in B$ беше произволно,
    \[
        \int_0^T p(t) dt \leq
        \inf_{(\tilde{x},\tilde{y})\in B} \lVert (\tilde{x}, \tilde{y}) - (x,\dot{x})\rVert =
        d((x,\dot{x}), B)
    \]
    в частност
    \[
        \int_0^T p(t) dt \leq
        \lVert (\bar{x}, \dot{\bar{x}}) - (x,\dot{x})\rVert \leq 1
    \]
    Значи $p\in L^1([0,T])$ и можем да приложим теорема \ref{Filippov} и да получим $(\tilde{x}, \dot{\tilde{x}})\in Q$, т.ч.
    \[
        \lVert x(t) - \tilde{x}(t) \rVert \leq \varphi(t)
    \]
    \[
        \lVert \dot{x}(t) - \dot{\tilde{x}}(t) \rVert \leq l(t) \varphi(t) + p(t)\ \text{п.н.}
    \]
    където
    \[
        \varphi(t) =
        \int_0^t p(s) e^{\int_s^t l(\tau) d\tau}ds \leq
        \int_0^t p(s) e^{\int_0^T l(\tau) d\tau}ds =
        e^L \int_0^t p(s) ds.
    \]
    С това можем да пресметнем
    \[
        \lVert x - \tilde{x} \rVert_\infty =
        \sup_{t\in [0,T]} \lVert x(t) - \tilde{x}(t)\rVert \leq
        \sup_{t\in [0,T]} e^L \int_0^t p(s) ds =
        e^L \int_0^T p(s) ds,
    \]
    \[
        \lVert \dot{x} - \dot{\tilde{x}} \rVert_1 =
        \int_0^T \lVert \dot{x}(s) - \dot{\tilde{x}}(s) \rVert ds \leq
        \int_0^T \left[l(s) e^L \int_0^s p(\tau)d\tau + p(s)\right] ds \leq
    \]
    \[
        e^L \int_0^T l(s) \int_0^T p(\tau)d\tau ds + \int_0^T p(s)ds =
        (L e^L + 1) \int_0^T p(s)ds
    \]
    и сумарно
    \[
        \lVert (x, \dot{x}) - (\tilde{x}, \dot{\tilde{x}}) \rVert =
        L \lVert x - \tilde{x} \rVert_\infty + \lVert \dot{x} - \dot{\tilde{x}} \rVert_1 \leq
        % L e^L \int_0^T p(s) ds + (L e^L + 1) \int_0^T p(s)ds =
        (2L e^L + 1)\int_0^T p(s)ds
    \]
    Както вече видяхме $\int_0^T p(t) dt \leq d((x,\dot{x}), B)$, така че
    \[
        d((x, \dot{x}), Q)\leq
        \lVert (x, \dot{x}) - (\tilde{x}, \dot{\tilde{x}}) \rVert \leq
        (2L e^L + 1) d((x,\dot{x}), B)
    \]
\end{proof}

Това, в комбинация с Теорема \ref{BigCone} и Твърдение \ref{TIP} ни дават желаната
\begin{theorem}
    \small
    \begin{displaymath}
        \Tangent_Q\Point \supseteq A \cap \{\Vec\in\Univ \mid \Vect\in T_{\Gr}(\Pointt)\ \text{п.н. в } [0,T]\}
    \end{displaymath}
\end{theorem}

